\section{Replica entropy and sparse isolated paths}\label{c:sec:replica}
Write $N=n=2^d$ in this section, and write $\lambda^\top$ for the
conjugate partition, so $\lambda^\top_1=\lambda'_1$.
The weak-entropy argument of Section~\ref{sec:c-baseline-forest}
used a true-sweep operator estimate after clipping. Here explicit
partner-repeat counts give the entropy input, and a separate
Hilbert--Schmidt estimate for isolated paths gives the completion.
The relative entropy of the image of a large random list need not
tend to zero: a bound
polynomial in $d$ suffices after truncating exceptionally concentrated
outputs. Representations with small dimension require a separate
argument, supplied here by cancellation of isolated paths in one sweep.
The averaging over input lists is essential throughout the entropy
calculation. For a law $w$ on $S_N$ and an ordered list $I$ of distinct input positions, write
$w_I$ for the law of its ordered image under a permutation sampled from $w$.
\begin{theorem}[Replica entropy route]\label{c:rep:main}
For the physical Thorp shuffle on $N=2^d$ cards,
\[
 \|q_d^{*(1327104d)}-U_{S_N}\|_{\TV}\longrightarrow0.
\]
More precisely, put $L=256$, $q=N-N/L$, and let $w$ be the law after $16384d$ physical
shuffles. Averaging over all ordered distinct $q$-tuples $I$, its projected law satisfies
\begin{equation}
 \E_I\KL(w_I\|\mathbf u_q)=O(d^4).\label{c:rep:e1}
\end{equation}
Here $\mathbf u_q$ is uniform on ordered distinct $q$-tuples, and
$\KL(p\|u)=\sum_xp(x)\log(p(x)/u(x))$, with zero summands interpreted as zero.
\end{theorem}
The replica below stays on its current baseline label except when it
encounters the tagged path. This differs from the walker of
Section~\ref{sec:c-baseline-forest}, which can change labels at any
active partner. The replica can return to the tagged label only when
that label's partner sequence repeats. Counting one repeat explicitly
and bounding two repeats by $O(d^4/N^2)$ yields the entropy estimate.
For the small-dimension estimate we reuse the deterministic path kernel
of Lemma~\ref{lem:c-sweep-kernel}, now in Hilbert--Schmidt norm. The
final assembly uses $80$ truncated blocks and one untruncated block.
All logarithms in this section are natural.

Fix constants and sizes as follows:
\[
L=256,\qquad r=N/L,\qquad q=N-r,\qquad A=64L,\qquad T=A d.
\]
Here \(r\) is an integer for all \(d\) under consideration. Put $G=S_N$ and write \(\kappa\)
for the law on
\(G\) of \(d\) shuffles, and \(w=\kappa^{*A}\), the law of \(T\) shuffles. We use the notation
\((n)_b=n(n-1)\cdots(n-b+1)\).

\subsection{Replica collisions and projection entropy}
We prove the entropy assertion
of Theorem~\ref{c:rep:main} by the chain rule, exposing the paths of one
card at a time.

\begin{proof}[Proof of the entropy bound]
\label{c:rep:entropy-proof}
Consider the chain-rule term for the \((j+1)\)-st card,
\(0\le j\le q-1\); the reference law from \(\mathbf u_q\), given the final positions of the
first \(j\) cards, is uniform on the \(M=N-j\) other positions. Further condition on
\(\Gamma\), the data of the whole trajectories through time \(T\) of these first \(j\) cards
(generated by the shuffles used for \(w\)). Call these \(j\) cards observed cards, and the other labels
unobserved. For the fixed list \(I\), write \(p_\Gamma\) for the law of the final
position of the next card, whose label we denote by \(X\), given \(\Gamma\). The positions
available to \(X\) at the end, complementary to the first \(j\) final positions, form the final
available set; write \(u_\Gamma\) for its uniform law.
Conditioning additionally as above gives an upper bound on the entropy term by convexity
of relative entropy, since the uniform reference here only depends on the first \(j\) final
positions. Thus the term, now averaged over \(I\), is bounded by
\begin{equation}
\mathbb E_{I,\Gamma}\mathrm{KL}(p_\Gamma\|u_\Gamma)
\ \le\ M\mathbb E_{I,\Gamma}\sum_z p_\Gamma(z)^2-1,                  \label{c:rep:e2}
\end{equation}
where \(z\) indexes the final available positions and the last inequality uses \(\log b\le b-1\).

By Lemma~\ref{lem:marginal-averaging}, conditional on \(I,\Gamma\),
switches visited by observed cards have fixed choices, while the coins
on pairs of available positions remain independent and fair. We now
construct the additional coupling needed to represent the squared mass.

\label{c:rep:replica-construction}
Generate a full shuffle trajectory (the reference trajectory for every label) using all actual
switch choices, so \(X\) uses its reference trajectory. We construct a replica path, starting
at the same position as \(X\), that will be on reference trajectories of unobserved labels. When at a
switch also being used by \(X\), and this switch has two available positions, the replica uses an extra
independent fair coin, choosing to follow either \(X\) or its partner in the reference
trajectories through this switch. In any other case, the replica follows the reference
trajectory of the label at its position through its switch.

The positional path of this replica and the path of \(X\) are independent copies given
\(I,\Gamma\). To see it, condition on the past at any stage in this construction. Given
\(I,\Gamma\), the single-path transition at an available position is forced if the switch contains an
observed card, and otherwise it uses the two outputs with equal probability. If the replica and
\(X\) are at different switches, any unforced choices there come from independent fresh coins.
If they use the same switch with only one available position, they both have a forced transition. If they use
the same switch with two available positions, the rule with the extra coin makes the replica transition
uniform independently of the transition of \(X\). Thus the joint positional transitions are
products at every step and the kernels in this description only require the respective present
positions given \(I,\Gamma\). In particular, writing \(C_s\) for the label tracked by the
replica at time \(s\) (\(C_0=X\)), we have
\[
\mathbb E_{I,\Gamma}\sum_z p_\Gamma(z)^2=\Pr(C_T=X),
\]
since the reference trajectories of the labels give distinct final positions.

We now average without conditioning on \(\Gamma\). Fix which label \(X\) is, and, in the
reference trajectories, let \(J_s\) be the label of the partner of \(X\) at step \(s\), where
the steps are indexed \(s=0,\ldots,T-1\). The set of \(M-1\) unobserved labels other than \(X\) in
this estimate is uniformly chosen among the other \(N-1\) labels, independently of the
reference trajectories, due to the choice of \(I\). Given the sequence \(J_s\) and the unobserved
labels, \(C_s\) can only change when it is in \(\{X,J_s\}\) and \(J_s\) is an unobserved label. In
this case \(C_{s+1}\) is \(X\) or \(J_s\) each with probability \(1/2\), by the extra coin.

\label{c:rep:repeat-proof}
We give estimates for repeats in the partner sequence, over the reference randomness. Write
\(b_{i,u}=\Pr(J_i=J_u)\). Then, for \(0\le i<u<T\),
\begin{equation}
\begin{cases}
 b_{i,u}=0 & 0<u-i<d,\\
 b_{i,u}=2/N & u-i=d,\\
 b_{i,u}\le (1+2/N)/N & u-i>d.
\end{cases}                                                        \label{c:rep:e3}
\end{equation}
Indeed, at step \(i\) the two actual labels \(X,J_i\) have opposite appended bits. These bits
preclude a repeat partner for a gap less than \(d\). At \(u\ge i+d\), a repeat requires the
appended bits for these two labels at step \(u-d\) to be opposite and those at steps
\(u-d+1,\ldots,u-1\) to match, and these conditions suffice. Given opposite bits at \(u-d\),
the next \(d-1\) steps have different switches for these labels, so the matching part has
probability \(2^{-(d-1)}\). When \(u-d=i\), we have the opposite bits for sure. When \(u-d>i\),
the bits there are opposite with probability one if the two are partners at \(u-d\), and
\(1/2\) otherwise, making \(b_{i,u}=(1+b_{i,u-d})/N\). For a preceding pair of times
\(i<v=u-d\), \(b_{i,v}\) is at most \(2/N\): a smaller gap than \(d\) gives zero, and for gap
at least \(d\), the \(d-1\) required bit matches each have conditional probability at most
\(1/2\) given the past. Here \(J_i\) is known by the times requiring these matches. This proves
\eqref{c:rep:e3}.

\label{c:rep:two-repeats}
Let \(R_0=T-|\{J_s:0\le s<T\}|\). We also have
\begin{equation}
\Pr(R_0\ge 2)=O(T^4/N^2).                                          \label{c:rep:e4}
\end{equation}
Indeed, fix two distinct repeat times $u,v$ and earlier indices
$i<u$, $h<v$, and ask for $J_u=J_i$ and $J_v=J_h$. By
\eqref{c:rep:e3}, both gaps must be at least $d$. Let $\mathcal F_s$
contain all reference switch coins at steps strictly before $s$.
The partner $J_s$ is $\mathcal F_s$-measurable. In particular $J_i$
is already known before the interval
$[u-d+1,u-1]$ of $d-1$ required bit matches, and the analogous
statement holds for $J_h$ and $[v-d+1,v-1]$.

If these intervals overlap, then $|u-v|<d$. On the joint repeat
event the two earlier labels must differ: otherwise the same label
would be a partner at both $u$ and $v$, contradicting
\eqref{c:rep:e3}. At an overlap step impose this inequality as part
of the event being bounded; it is measurable before that step.
The three distinct labels $X,J_i,J_h$ must append the same bit.
If any two occupy one switch the event is impossible, because the
two true outputs are opposite. Otherwise three independent fair
switch coins give the event probability $1/4$. At a step in just one
matching interval the probability is at most $1/2$. Iterating these
bounds in chronological order gives $2^{-2(d-1)}=4/N^2$ for the fixed
quadruple $(i,u,h,v)$. The union bound over $O(T^4)$ quadruples proves
\eqref{c:rep:e4}.

We have controlled the frequency of repeated partners. It remains to
average which distinct partners belong to the hidden set. This step is
where the uniform random input list enters.

Let
\[
\delta=\frac{M-1}{N-1},\qquad \eta=1-\frac{\delta}{2}.
\]
Here \(\delta\ge 1/L\). Given \(J_s\), mark each distinct label in the
sequence according to whether it is unobserved. The probabilities of these
mark patterns are upper bounded by those for independent marks of chance
\(\delta\), times a common factor
\[
c_T=\frac{(N-1)^T}{(N-1)_T}=1+O(T^2/N)
\]
for large \(d\). In fact for \(k\) distinct labels with \(b\) specified to be unobserved labels and
\(k-b\) to be observed, the true probability is \((M-1)_b(N-M)_{k-b}/(N-1)_k\), and \(k\le
T=o(\sqrt N)\).

Suppose \(R_0\le 1\). The independent-marks bound (before the factor \(c_T\)) for never leaving
\(X\), i.e. \(C_s=X\) throughout, is at most \(\eta^{T-1}\), using a step with each of at least
\(T-1\) distinct partners. If \(C_s\) does leave \(X\) and \(C_T=X\), there must be a repeat
\(J_i=J_u\), which for \(R_0\le 1\) is the only one. At \(i\) it must switch from \(X\) to the
repeated label, and at \(u\) back to \(X\). There can be no other departure. The
independent-marks probability for this behavior given such a sequence is
\[
\frac{\delta}{4}\,\eta^{i+T-1-u}.
\]
Here we used the mark of the repeated label and two particular coin results there; all the
steps requiring no departure before \(i\) or after \(u\) have different partners, also
different from the repeated one.

Combining this with \eqref{c:rep:e3} and \eqref{c:rep:e4},
\[
\begin{aligned}
\Pr(C_T=X)
 &\le O(T^4/N^2)+c_T\left[\eta^{T-1}
       +\frac{\delta}{4}\sum_{i<u}\eta^{i+T-1-u} b_{i,u}\right]\\
 &\le O(T^4/N^2)+c_T\left[\eta^{T-1}
       +\frac{1+2/N}{\delta N}
       +\frac{\delta}{4N}\,T\eta^{T-1-d}\right].
\end{aligned}
\]
Indeed the non-special bound \((1+2/N)/N\) gives the middle term by a double geometric sum
since \(\sum_{l\ge 0}\eta^l=2/\delta\), and an additional bound of \(1/N\) can be included for
gaps \(d\). We have \(\eta^{T-1-d}\le N^{-4}\) for large \(d\), by \(A=64L\), and \(M/(\delta
N)=1+O(1/N)\). Consequently
\[
M\Pr(C_T=X)-1\ \le\ O(T^4/N)
\]
uniformly in \(j\), under the averaging used in \eqref{c:rep:e2}.
Each conditional entropy term is therefore $O(T^4/N)$. There are
$q\le N$ terms, so their sum is $O(T^4)=O(d^4)$, proving
\eqref{c:rep:e1}.
\end{proof}

\subsection{Applying the entropy cutoff}
\label{c:rep:trimming}
Choose a uniform partition of the input positions into $L=256$ equal
blocks. Each complementary list is a uniform $q$-subset; adding a
uniform order gives the list law in \eqref{c:rep:e1}. Reordering a list
does not change the entropy of its image. Averaging therefore supplies
one deterministic partition with complementary lists $I_1,\ldots,I_L$
satisfying
\[
 \sum_{a=1}^{L}\mathrm{KL}(w_{I_a}\|\mathbf u_q)=O(d^4).
\]
Write $H_a$ for the permutations supported on block $a$, so
$|H_a|/|G|=1/(N)_q$.

We apply the entropy-truncation and isotypic transfer of
\cite[Theorem~\ref*{l-l:replica-clipping}]{partialFourier}.
Take $L=256$ blocks, total relative entropy $H=O(d^4)$,
logarithmic cutoff $b=d^6$, and reciprocal-degree exponent $u=32$. The theorem
restricts $w$ to the event that every projected density is at most
$e^b$, giving a subprobability $v\le w$ of mass $m$ such that
\begin{equation}\label{c:rep:mass}
 1-m\le\frac{H+L/e}{b}=O(d^{-2}).
\end{equation}
For clarity, this is one restriction of the full law; the $L$ marginal
laws are not clipped independently. It yields simultaneously
\begin{equation}\label{c:rep:e5}
 \sum_{h\in H_a}(\check v*v)(h)
 \le e^{d^6}\frac{|H_a|}{|G|},\qquad
 \check v(g)=v(g^{-1}),
\end{equation}
and
\begin{equation}\label{c:rep:e8}
 \|\widehat v(\lambda)\|_{\mathrm{HS}}^2
 \le LC e^{d^6}D_\lambda^{-1+34/L}.
\end{equation}
Here and below the Hilbert--Schmidt norm uses ordinary, unnormalized
trace. Here $34/L=(32+2)/L$: the reciprocal-degree bound contributes $32/L$,
the squared subgroup dimension contributes $2/L$, and the ambient dimension
contributes $-1$. No fixed-list total-variation
estimate has been used.

We also use the degree bounds
\begin{equation}\label{c:rep:e6}
 \log D_\lambda\ge\frac{\log2}{2}\sqrt{\ell(\lambda)},
 \qquad
 \sup_{s\ge1}\sum_{\alpha\vdash s}D_\alpha^{-32}\le C,
 \qquad
 \ell(\lambda)=N-\max(\lambda_1,\lambda'_1),
\end{equation}
from \cite[Lemma~\ref*{l-l:degree-sqrt-deficit}]{partialFourier}.
Here $\lambda'_1$ is the first-column length. Since
$1-34/256>1/2$, \eqref{c:rep:e8} gives
\begin{equation}\label{c:rep:large-dimension}
 \|\widehat v(\lambda)\|_{\mathrm{HS}}^2\le D_\lambda^{-1/2}
 \qquad(\log D_\lambda>d^8)
\end{equation}
for all sufficiently large $d$: the fixed factor $LC$ and the cutoff
cost $d^6$ are smaller than
$(1/2-34/256)\log D_\lambda$ in this range.

\subsection{An untruncated sweep on the remaining representations}
The cutoff has controlled every representation with
$\log D_\lambda>d^8$. We now use the true sweep law $\kappa$ on the
remaining shapes. This factor will be kept untruncated in the final
convolution, so that its exact cancellations remain available.

For the remaining estimates consider \(\log D_\lambda\le d^8\), now using the untruncated law.
By \eqref{c:rep:e6},
\[
N-\max(\lambda_1,\lambda^\top_1)\ \le\ O(d^{16}).
\]
Those with \(\lambda^\top_1=N-k\), \(k\le O(d^{16})\), have \(\widehat\kappa(\lambda)=0\) for
large \(d\). Indeed, the first shuffle averages over the subgroup generated by the \(N/2\)
disjoint flips of partner input pairs, and includes a deterministic string rotation
(\((s,z)\mapsto(z,s)\) for bit \(s\) and suffix \(z\)) after this. The average for the flip
subgroup is a projection onto invariants. In \(\lambda\), the sign twist of \(\lambda^\top\),
there can be no such invariants. To see this, place \(\lambda^\top\) inside the representation
on ordered distinct \(k\)-tuples as allowed by branching. A vector that would be invariant upon
sign twisting has to transform by sign under all the pair flips. For each tuple there is a pair
unused by it as \(k<N/2\); its flip fixes the tuple, so the entry of the vector there must
vanish. This proves the asserted absence, and hence the vanishing for the law including the
first shuffle.

\label{c:rep:sparse-proof}
\begin{proposition}[Sparse isolated-path estimate]
\label{c:rep:sparse}
Fix $C_0<\infty$. If $\lambda\vdash N=2^d$ has first-row length
$N-k$, where $1\le k\le C_0d^{16}$, then, for all sufficiently large
$d$ depending only on $C_0$, the law $\kappa$ of one sweep satisfies
\begin{equation}
\|\widehat\kappa(\lambda)\|_{\mathrm{HS}}\le N^{-k/10}.                \label{c:rep:e9}
\end{equation}
\end{proposition}
\begin{proof}
Use the representation on ordered distinct \(k\)-tuples, with a group element taking the
coordinate vector for \(\boldsymbol x\) to that for \(g\boldsymbol x\).
The law $\kappa=q_d^{*d}$ is the true-sweep law $K$ of
Lemma~\ref{lem:c-sweep-kernel}: after $d$ physical shuffles the moving
frame returns to the identity. For ordered distinct
\(\boldsymbol x,\boldsymbol y\), use that lemma's kernels
\[
F_S(\boldsymbol x,\boldsymbol y)=N^{|S|}\Pr_\kappa(g x_a=y_a\text{ for }a\in S),
\qquad
F_\circ=F_*=\sum_{S\subseteq[k]}(-1)^{k-|S|}F_S.
\]
The averaged matrix for \(\kappa\) has entries (output
\(\boldsymbol y\), input \(\boldsymbol x\)) given by \(F_{\{1,\ldots,k\}}/N^k\). When
projecting on the output side to \(\lambda\)-type, we can replace \(F_{\{1,\ldots,k\}}\) by
\(F_\circ\): the other terms map into functions of proper subsets of output positions only. The
space for any such subset is a copy of a tuple module on fewer entries, not containing
\(\lambda\) by branching. Also \(\lambda\) does occur in our \(k\)-tuple module, and the
projected \(\kappa\)-matrix includes at least one copy of \(\widehat\kappa(\lambda)\).
Consequently, by taking the Hilbert-Schmidt bound before applying the output projection to the
replacement matrix,
\begin{equation}
\|\widehat\kappa(\lambda)\|_{\mathrm{HS}}^2
\le \left(\frac{(N)_k}{N^k}\right)^2 \mathbb E_{\boldsymbol x,\boldsymbol y}
                       \left|F_\circ(\boldsymbol x,\boldsymbol y)\right|^2,   \label{c:rep:e10}
\end{equation}
where each of \(\boldsymbol x,\boldsymbol y\) is separately a uniform tuple with distinct
positions and they are independent.

By Lemma~\ref{lem:c-sweep-kernel}, $|F_\circ|\le2^kk^{k/2}$,
and $F_\circ$ vanishes when the graph of shared prescribed switches
has an isolated vertex. Unlike the operator estimate in
Proposition~\ref{prop:c-true-sweep-sparse}, \eqref{c:rep:e10} requires
the probability of the remaining pairs of tuples, rather than a
bound on their bilinear form.
For \(k=1\), \(F_\circ\) already vanishes by isolation. For \(k\ge 2\), if there are no
isolated vertices then the graph contains a spanning forest with no isolated vertices, having
at least \(k/2\) and at most \(k-1\) edges. Consider first drawing all positions of the two
tuples independently with replacement. A shared switch in direction $i$
requires agreement of the length-$(d-i)$ input suffixes and the
length-$(i-1)$ output prefixes. Given the data at one vertex, a fresh vertex forms an
edge to it with probability at most \(2d/N\), by summing over the steps the probabilities of
suffix and prefix agreement. Thus for any particular forest the probability that all its edges
are present is at most \((2d/N)^{\#\text{edges}}\), by drawing along the trees. Conditioning
each tuple to have distinct positions costs a factor at most 2 for large \(d\) in our size
range. Summing over forests (at most \(k^{2k+1}\), by counting edge lists) bounds the
probability in \eqref{c:rep:e10} that \(F_\circ\ne 0\) by
\[
2 k^{2k+1}(2d/N)^{k/2}.
\]
Together with the size bound this yields \eqref{c:rep:e9}, since factors \(k^{O(k)}d^{O(k)}\)
are only \(\exp(O(k\log d))\) here, compared to \(N^{k/2}\) in the probability denominator.
\end{proof}

\subsection{Assembly with one untruncated factor} \label{c:rep:assembly}
\begin{proof}[Completion of Theorem~\ref{c:rep:main}]
Let \(B=80\) and \(\theta=v^{*B}* w\), of mass \(m^B\) and dominated pointwise by the true law
\(w^{*(B+1)}\). For \(\log D_\lambda>d^8\), \eqref{c:rep:e8} gives
\[
\|\widehat\theta(\lambda)\|_{\mathrm{HS}}^2\le D_\lambda^{-B/2}
\]
for large \(d\), by convolution multiplication and Hilbert-Schmidt submultiplicativity, using
also \(\|\widehat w(\lambda)\|_{\mathrm{op}}\le 1\) since it averages unitaries. It follows
from \eqref{c:rep:e6} that
\[
\sum_{\lambda:\log D_\lambda>d^8}D_\lambda \|\widehat\theta(\lambda)\|_{\mathrm{HS}}^2
\le C \exp\big((33-B/2)d^8\big)=o(1).
\]
At smaller dimensions, excluding the trivial representation, we have either the vanishing for
\(\kappa\) above (hence for \(\theta\)), or the bounds \eqref{c:rep:e9} with \(k\ge 1\), and
\(\|\widehat v(\lambda)\|_{\mathrm{op}}\le 1\). In the latter case \(D_\lambda\le N^k\) since
the \(k\)-tuple representation contains \(\lambda\), and \(w=\kappa^{*A}\), so
\[
D_\lambda \|\widehat\theta(\lambda)\|_{\mathrm{HS}}^2
\le N^k\big(N^{-k/10}\big)^{2A}.
\]
Summing this over these diagrams gives \(o(1)\) (the count for first row \(N-k\) is at most
\(\exp(3\sqrt k)\)).

Writing \(U_G\) for uniform on \(G\), Cauchy-Schwarz and Plancherel give
\[
\|\theta-m^B U_G\|_1^2
\le \sum_{\lambda\ \mathrm{nontrivial}}D_\lambda\|\widehat\theta(\lambda)\|_{\mathrm{HS}}^2
=o(1),
\]
with \(\ell^1\) on group masses. Hence TV distance for \(w^{*(B+1)}\) from uniform is bounded
by \(1-m^B+\tfrac12\|\theta-m^B U_G\|_1\), tending to zero. This law uses \((B+1)A d\) complete
shuffles, proving the asserted time $81\cdot16384d=1327104d$. Translation of the initial deck
preserves the distance to uniform, so the convergence is uniform over starting decks.
\end{proof}
